\subsection{Ontology-Based Information Extraction for ESG ABSA}

\subsubsection{Ontology-based IE as a structured alternative to bag-of-words approaches}

Conceptual shift and technical rationale. Ontology-based information extraction (IE) provides a formal, machine-readable representation of domain knowledge that encapsulates concepts, relations, and constraints. Unlike bag-of-words or flat keyword approaches, an ontology enables explicit modeling of ESG concepts (e.g., emissions, governance oversight, workforce diversity) and the relationships among them (hierarchies, part-of, associated indicators). This structure supports disambiguation, facet-level reasoning, and regulator-aligned querying, which are essential for credible ABSA in dense sustainability disclosures \cite{10.2139/ssrn.3738629,10.20944/preprints202407.2036.v1,10.7849/ksnre.2022.0010,10.47473/2020rmm0103}
Practical architecture. An ontology-driven IE pipeline typically comprises (a) a curated ESG ontology aligned to GRI/SASB/TCFD concepts, (b) a gazetteer of bilingual terms and synonyms for Indonesian/English, (c) a semantic layer that maps textual spans to ontology nodes via entity recognition and relation extraction, and (d) a reasoning component that derives higher-order signals (e.g., “emissions reduced” linked to an environmental KPI) from mapped facts. This contrasts with pure bag-of-words signals, which lack structural grounding and are vulnerable to lexical drift across languages and regulatory lexicons \cite{10.7849/ksnre.2022.0010,10.47473/2020rmm0103}

Novelty for ESG ABSA. Integrating ontology-driven IE with ABSA enables sentence- and clause-level sentiment to be tied not only to a surface term but to an ontology node with provenance, regulatory anchors, and hierarchical context. This supports explainable ABSA outputs where a sentiment claim is anchored to a regulatory topic (GRI/SASB/TCFD) and to a specific evidence path within the document, thereby enhancing auditability and regulatory traceability \cite{10.20944/preprints202407.2036.v1,10.47473/2020rmm0103}
\subsubsection{ESG taxonomies used in practice: GRI, SASB, and TCFD, and their integration}

GRI as a stakeholder-oriented taxonomy. GRI provides broad, multi-stakeholder content covering environmental, social, and governance topics, with materiality oriented toward sustainability impacts and stakeholder engagement. In ontology terms, GRI concepts anchor high-level domains and provide cross-cutting topics that populate the upper layers of the ESG ontology. This broad coverage supports signaling about governance and social impact beyond purely financial metrics, but its indicators are often qualitative and narrative-driven, necessitating structured extraction to achieve comparability \cite{10.1177/09726225241237304,10.1111/ejed.70225}.
SASB as a financial-materiality taxonomy. SASB emphasizes financially material ESG topics by industry, guiding disclosures that matter for investors and financial decision-making. In an ontology, SASB categories map to more granular sub-classes and indicators with explicit KPI targets, enabling ABSA to anchor sentiment to financially salient dimensions and to connect qualitative statements with quantitative performance data. The SASB focus complements GRI by increasing comparability across firms from an investment lens \cite{10.1177/09726225241237304,10.1111/ejed.70225}.
TCFD as a climate-risk framework. TCFD provides a structured framework around governance, strategy, risk management, and metrics/targets related to climate risk. Ontology mappings can encode these four pillars and their sub-indicators, enabling sentence-level sentiment to be interpreted in the context of climate-related disclosures, scenario analysis, and transition plans. This alignment is critical for ABSA when climate-related language co-occurs with other ESG topics \cite{10.1002/csr.70116,10.1017/s135732172300017x}
Cross-framework integration. The literature documents ongoing efforts to harmonize GRI, SASB, TCFD, and related frameworks into integrated reporting, often via mapping tables and dialog between standard-setters. For ABSA, integrating these frameworks within a single ontology enables cross-framework topic labeling, facilitates comparability across jurisdictions, and supports regulatory analysis in Indonesia and globally \cite{10.1007/978-3-319-04723-2_7,10.1108/ijlma-06-2022-0113,10.1145/3503044}
Mapping bilingual language to ontology nodes. A central practical challenge is aligning Indonesian and English report language to the formal ontology. This requires bilingual term mappings, cross-language entity normalization, and alignment of regulatory lexicons (Indonesian OJK/IDX terminology with international ESG terms). Ontology-based IE with bilingual gazetteers and cross-lingual embeddings can bridge language gaps and preserve semantic fidelity across languages \cite{10.1016/j.inffus.2025.103073,10.1111/eufm.12509,10.48550/arxiv.2511.00024}
\subsubsection{How ontology paths guide aspect assignment in automated ESG analysis}

Ontology-driven aspect assignment. In ABSA, each extracted sentiment instance is associated with an ESG aspect. An ontology path provides the semantic trajectory: from text mention to a concept node (e.g., “emissions intensity” → Environmental → Emissions) and further to sub-aspects or indicators (e.g., Scope 1/2 emissions, energy efficiency). Path-based reasoning supports disambiguation when surface terms are ambiguous or multi-annotated within the same sentence, enabling precise topic assignment and richer downstream analyses \cite{10.2139/ssrn.3738629,10.7849/ksnre.2022.0010}
Reasoning and constraint mechanisms. Ontologies enable logical inferences, such as recognizing that a sentence about “net-zero targets” is climate-related and should be linked to the TCFD/TCFD-aligned metrics, even if the text does not explicitly mention “TCFD.” Constraints (domain axioms) ensure consistent labeling across languages and ensure that sentiment about an environmental target maps to the corresponding KPI or regulatory anchor. This fosters consistency in ABSA outputs under multilingual and multi-framework contexts \cite{10.1177/09726225241237304,10.1111/ejed.70225,10.1145/3503044}
Explainability through ontology traces. By exposing the path from a surface phrase to ontology nodes and then to regulatory anchors, the system can generate explanations such as: “The sentiment toward emissions is positive because the sentence maps to Environmental > Emissions > Target KPIs per SASB, with evidence in the KPI table.” Such provenance traces satisfy regulator expectations for auditability and support user trust in automated ABSA outputs \cite{10.1111/reel.12554,10.1111/1475-679x.12575,10.3390/su11041093}
Handling multi-aspect sentences. OMS (Ontological Multi-Shift) paths allow a single sentence to map to multiple aspects if required (e.g., a sentence could discuss both emissions and supply-chain governance). The ontology provides disambiguation cues, such as co-occurring indicators and regulatory tags, enabling multi-label aspect assignment while preserving interpretability \cite{10.3390/su11041093,10.7849/ksnre.2022.0010}
.
\subsubsection{Knowledge graph applications for ESG}

Knowledge graphs for ESG governance. ESG knowledge graphs integrate entities (companies, frameworks, indicators, targets), relations (owns, reportsto, measures, targetof), and evidence (sources, pages, sections). They enable semantically rich querying (e.g., “which companies report emissions reductions within SASB material topics for the energy sector?”) and support complex reasoning over regulatory alignment, materiality, and performance signals. This approach extends ABSA by providing a flexible, queryable backbone for cross-document, cross-language analytics \cite{10.1145/3503044}

Use cases in ABSA workflows. In a practical ABSA pipeline, a knowledge graph can: (a) store ontology-driven mappings from sentences to ESG concepts, (b) integrate numerical KPIs and targets for outcomes, (c) support explainable ABSA by enabling path-based explanations, and (d) enable cross-document linking (e.g., linking a governance claim in a narrative to a governance structure in the annual report). Graph-based reasoning supports more robust evidence retrieval and regulator-aligned reporting \cite{10.1145/3503044}
Cross-language knowledge graphs. For bilingual Indonesian/English reports, multilingual knowledge graphs leverage cross-lingual alignment to map terms to a shared ontology, enabling consistent interpretation across languages and ensuring that regulatory anchors remain coherent when language switches occur within a document or across documents \cite{10.1016/j.inffus.2025.103073,10.1111/eufm.12509,10.48550/arxiv.2511.00024}
Practical benefits and evaluation. Knowledge graphs improve traceability, support dynamic updates as ESG frameworks evolve, and enable scalable queries across large corpora of reports. Evaluation focuses on correctness of entity and relation extraction, completeness of ontology-path mappings, and usefulness of graph-based queries for regulators and investors \cite{10.1145/3503044}
Notes on challenges and synthesis

Language-to-ontology mapping is the principal technical bottleneck. The bilingual Indonesian/English setting introduces lexical variability, code-switching, and region-specific regulatory lexicons, requiring robust cross-lingual alignment, lexicon curation, and continuous maintenance of multilingual synonyms. Approaches include multilingual embeddings, translation-based alignment, and ontology-informed normalization procedures to mitigate translation-induced semantic drift \cite{10.1016/j.inffus.2025.103073,10.1111/eufm.12509,10.48550/arxiv.2511.00024}
Evaluation and governance. Ontology-based IE requires specialized evaluation protocols: concept coverage, relation extraction accuracy, path consistency, and regulatory alignment fidelity. It also requires alignment with existing standards (GRI, SASB, TCFD) to ensure practitioner acceptance and regulatory relevance \cite{10.1177/09726225241237304,10.1111/ejed.70225}
Summary. Ontology-based IE provides a principled, scalable, and auditable route to ESG ABSA, moving beyond bag-of-words to structured semantic reasoning. By anchoring aspects to robust ESG taxonomies, guiding ABSA through ontology paths, and leveraging knowledge graphs for evidence retrieval and governance, this approach enhances interpretability, cross-language consistency, and regulatory alignment in Indonesian and international ESG disclosure analysis \cite{10.2139/ssrn.3738629,10.7849/ksnre.2022.0010,10.1145/3503044}

% Recommended references for further grounding (IEEE style)


% 110–112 OECD and standard-setting literature on taxonomy harmonization and regulatory alignment (various sources cited in the broad ESG standards discourse).
% End of section.